\documentclass[12pt,a4paper]{article}

% ------------------------------------------------
% Pacotes
% ------------------------------------------------
\usepackage[utf8]{inputenc}
\usepackage[T1]{fontenc}
\usepackage[brazil]{babel}
\usepackage{graphicx}
\usepackage[table]{xcolor}
\usepackage{array}
\usepackage{tabularx}
\usepackage{geometry}
\usepackage{fancyhdr}
\usepackage{mathptmx} % Times New Roman
\usepackage{enumitem}\usepackage{amsfonts}
\usepackage{amsmath} % Recomendado para o comando \text{}
\usepackage{amssymb}
\usepackage{tikz}
\usetikzlibrary{positioning, shapes.geometric, arrows.meta}

% TikZ Styles for ER and Logical Diagrams
\tikzset{
    block/.style={rectangle, draw, thick, minimum width=2.5cm, minimum height=1cm, fill=white, align=center},
    diamond_node/.style={diamond, draw, thick, minimum width=2cm, minimum height=1cm, fill=white, aspect=2},
    attribute/.style={ellipse, draw, thin, minimum width=1.5cm, minimum height=0.6cm, fill=white, font=\scriptsize},
    line/.style={draw, thick},
    table_node/.style={rectangle, draw, thick, minimum width=3cm, minimum height=0.8cm, fill=gray!10, font=\bfseries, align=center},
    column_node/.style={rectangle, draw, thin, minimum width=3cm, minimum height=0.6cm, fill=white, font=\small, align=left}
}

% ------------------------------------------------ 
% Configuração de página
% ------------------------------------------------
\geometry{
    a4paper,
    left=2cm,
    right=2cm,
    top=2cm,
    bottom=2cm,
    headheight=15pt
}

% Caixa para resposta discursiva
\newcommand{\resposta}[1]{
\noindent
\framebox[\textwidth]{
\begin{minipage}{0.98\textwidth}
\vspace{#1}
\end{minipage}
}
}

\pagestyle{fancy}
\fancyhf{}
\renewcommand{\headrulewidth}{0pt}

\setlength{\parindent}{0pt}

% ------------------------------------------------
% Comandos auxiliares
% ------------------------------------------------
\newcolumntype{L}[1]{>{\raggedright\arraybackslash}p{#1}}
\newcolumntype{C}[1]{>{\centering\arraybackslash}p{#1}}

% ------------------------------------------------
% Documento
% ------------------------------------------------
\begin{document}

% ------------------------------------------------
% Cabeçalho
% ------------------------------------------------
\begin{tabularx}{\textwidth}{@{} L{4cm} X @{}}
    \raisebox{-0.5cm}{\includegraphics[width=4cm]{Universidade_de_Rio_Verde.jpg}} &
    \hfill
    \begin{tabular}{|>{\centering\arraybackslash}m{10cm}|>{\raggedright\arraybackslash}p{1.8cm}|}
        \hline
        \cellcolor{black}\color{white}\rule{0pt}{1.2cm}\textbf{AVALIAÇÃO 2 - SEGUNDA CHAMADA} 
        & 
        \scriptsize Nota \tabularnewline
        \hline
    \end{tabular} 
\end{tabularx}
\vspace{0.3cm}

% ------------------------------------------------  
% Tabela principal
% ------------------------------------------------
\begin{tabularx}{\textwidth}{|X|L{2cm}|}
\hline
\textbf{Curso} & \textbf{Campus} \\
\textit{Engenharia de Software} & \textit{Rio Verde} \\
\hline
\multicolumn{2}{|l|}{\textbf{Disciplina}} \\  
\multicolumn{2}{|l|}{ \textit{Banco de dados}} \\
\hline
\textbf{Nome do(a) acadêmico(a)} & \textbf{Valor} \\
ISABELLE CRUVINEL SANTOS OLIVEIRA & 7,0 \\
\hline
\textbf{Professor} & \textbf{Turma} \\
\textit{ Sergio Souza Novak} & \textit{B} \\
\hline
\textbf{Data da Avaliação} & \textbf{Código} \\
\textit{18/05/2026} & \textit{ESW417} \\
\hline
\end{tabularx}
 
\vspace{0.2cm}

\begin{tabularx}{\textwidth}{|X|}
\hline
\cellcolor{black}\color{white}
\centering \textbf{ATENÇÃO:} Somente serão passíveis de \textbf{REVISÃO} avaliações resolvidas à \textbf{TINTA}. \tabularnewline
\hline
\end{tabularx}

\vspace{0.3cm}

\textbf{ORIENTAÇÕES PARA A RESOLUÇÃO}  
\begin{itemize}
    \item A avaliação é individual, sem consulta;
    \item A interpretação de texto faz parte da avaliação;
\end{itemize}

\vspace{0.5cm}

\begin{center}
    \rule{\textwidth}{0.4pt}
    \textbf{QUESTÕES DISCURSIVAS}
\end{center}

\vspace{0.2cm}

\textbf{INSTRUÇÕES:} Esta seção contém \textbf{10 exercícios discursivos}. Recomenda-se reservar \textbf{100 minutos} para sua resolução.

\vspace{0.3cm}

\begin{enumerate}
    \item \textbf{(BCNF)} Analise o cenário de orientação acadêmica abaixo:
    \begin{center}
        \begin{tikzpicture}[node distance=1.5cm, every node/.style={font=\small}]
            \node[block] (aluno) {Aluno};
            \node[block, right=5cm of aluno] (prof) {Professor};
            \node[diamond_node, right=1.2cm of aluno] (orient) {Orientação};
            \node[attribute, above=0.6cm of orient] (disc) {Disciplina};
            \node[attribute, above=0.6cm of prof] (email) {Email};
            \draw[line] (aluno) -- (orient); \draw[line] (prof) -- (orient);
            \draw[line] (orient) -- (disc); \draw[line] (prof) -- (email);
        \end{tikzpicture}
    \end{center}
    \textbf{Modelo Lógico:} \texttt{ORIENTACAO (Aluno, Professor, Disciplina, Email)} \\
    \textbf{Regras:} Cada professor tem um \textbf{Email} único e leciona apenas uma \textbf{Disciplina}. Para cada disciplina, o aluno tem um único orientador. Explique a violação da BCNF e normalize.
    \resposta{5cm}

    \item Realize a normalização para a 1FN da tabela abaixo, explicando o motivo técnico:
    \begin{center}
        \small
        \begin{tabular}{|c|c|L{7cm}|}
            \hline
            \cellcolor{gray!20}\textbf{ID} & \cellcolor{gray!20}\textbf{Funcionário} & \cellcolor{gray!20}\textbf{Dependentes} \\ \hline
            101 & João Silva & Maria (Filha), Pedro (Filho) \\ \hline
            102 & Fernanda Lima & Lucas (Filho) \\ \hline
        \end{tabular}
    \end{center}
    \resposta{4cm}

    \item Analise o modelo lógico da tabela \textbf{PEDIDOS} abaixo e identifique as dependências funcionais existentes (parciais ou transitivas). Em seguida, realize a decomposição necessária para que o modelo atinja a \textbf{Terceira Forma Normal (3FN)}:
    \begin{center}
        \begin{tikzpicture}[node distance=0cm, outer sep=0pt]
            \node[table_node] (ped) {PEDIDOS};
            \node[column_node, below=of ped] (c1) {\underline{ID\_Pedido}};
            \node[column_node, below=0.6cm of ped] (c2) {Data};
            \node[column_node, below=1.2cm of ped] (c3) {Cod\_Livro};
            \node[column_node, below=1.8cm of ped] (c4) {Titulo};
            \node[column_node, below=2.4cm of ped] (c5) {Preco};
        \end{tikzpicture}
    \end{center}
    \resposta{4cm}

    \item Analise o modelo lógico da tabela \textbf{ACADEMICO} abaixo. Explique o conceito de Dependência de Valores Múltiplos (MVD) e realize a decomposição necessária para que o modelo atinja a \textbf{Quarta Forma Normal (4FN)}, considerando que as disciplinas e os projetos são independentes:
    \begin{center}
        \begin{tikzpicture}[node distance=0cm, outer sep=0pt]
            \node[table_node] (acad) {ACADEMICO};
            \node[column_node, below=of acad] (a1) {\underline{Professor}};
            \node[column_node, below=0.6cm of acad] (a2) {\underline{Disciplina}};
            \node[column_node, below=1.2cm of acad] (a3) {\underline{Projeto\_Pesquisa}};
        \end{tikzpicture}
    \end{center}
    \resposta{4cm}

    \item Para cada diagrama de Venn, identifique o \textbf{JOIN} e escreva um parágrafo relacionando a figura ao comando SQL correspondente entre as tabelas \textbf{A} e \textbf{B} do modelo abaixo:\\
    \begin{center}
        \begin{tikzpicture}[scale=0.5]
            \begin{scope} \clip (0,0) circle (1cm); \fill[gray!40] (1.2,0) circle (1cm); \end{scope}
            \draw (0,0) circle (1cm) node {A}; \draw (1.2,0) circle (1cm) node {B}; \node at (0.6, -1.5) {Fig I};
        \end{tikzpicture} \hspace{0.5cm}
        \begin{tikzpicture}[scale=0.5]
            \fill[gray!40] (0,0) circle (1cm); \draw (0,0) circle (1cm) node {A}; \draw (1.2,0) circle (1cm) node {B}; \node at (0.6, -1.5) {Fig II};
        \end{tikzpicture} \hspace{0.5cm}
        \begin{tikzpicture}[scale=0.5]
            \fill[gray!40] (1.2,0) circle (1cm); \draw (0,0) circle (1cm) node {A}; \draw (1.2,0) circle (1cm) node {B}; \node at (0.6, -1.5) {Fig III};
        \end{tikzpicture} \hspace{0.5cm}
        \begin{tikzpicture}[scale=0.5]
            \fill[gray!40] (0,0) circle (1cm); \fill[gray!40] (1.2,0) circle (1cm); \draw (0,0) circle (1cm) node {A}; \draw (1.2,0) circle (1cm) node {B}; \node at (0.6, -1.5) {Fig IV};
        \end{tikzpicture}
    \end{center}
    \begin{center}
        \begin{tikzpicture}[node distance=0cm]
            \node[table_node] (tabA) {Tabela A};
            \node[column_node, below=of tabA] {id\_a (PK)};
            \node[column_node, below=0.6cm of tabA] {nome\_a};
            \node[table_node, right=2cm of tabA] (tabB) {Tabela B};
            \node[column_node, below=of tabB] {id\_b (PK)};
            \node[column_node, below=0.6cm of tabB] {fk\_a (FK)};
        \end{tikzpicture}
    \end{center}
    \resposta{4cm}

    \newpage

    \textbf{CONTEXTO PARA AS QUESTÕES 6 A 10:} Considere o modelo lógico de um sistema de vendas abaixo:
    \begin{center}
        \begin{tikzpicture}[node distance=0cm, outer sep=0pt]
            \node[table_node] (prod) {Produtos};
            \node[column_node, below=of prod] (p1) {\underline{id\_prod}};
            \node[column_node, below=0.6cm of prod] (p2) {nome\_produto};
            \node[column_node, below=1.2cm of prod] (p3) {categoria};
            \node[column_node, below=1.8cm of prod] (p4) {preco};
            
            \node[table_node, right=3cm of prod] (ped) {Pedidos};
            \node[column_node, below=of ped] (pe1) {\underline{id\_pedido}};
            \node[column_node, below=0.6cm of ped] (pe2) {data\_pedido};
            \node[column_node, below=1.2cm of ped] (pe3) {quantidade};
            \node[column_node, below=1.8cm of ped] (pe4) {fk\_prod (FK)};
        \end{tikzpicture}
    \end{center}

    \item Escreva o comando SQL (\textbf{INSERT}) para adicionar um novo produto na tabela \texttt{Produtos}. O produto é um 'Monitor Gamer' da categoria 'Eletrônicos' com o preço de R\$ 1.200,00 (Assuma \texttt{id\_prod} = 50).\\
    \resposta{3cm}

    \item Elabore uma consulta SQL (\textbf{SELECT}) que liste o nome e o preço de todos os produtos da categoria 'Eletrônicos', ordenando os resultados do maior para o menor preço (\textbf{ORDER BY}).\\
    \resposta{3cm}

    \item Qual consulta SQL utiliza corretamente a função de agregação para calcular a média de preços de todos os produtos cadastrados na tabela \texttt{Produtos}?
    \begin{enumerate}[label=\Alph*)]
        \item \texttt{SELECT MEAN(preco) FROM Produtos;}
        \item \texttt{SELECT AVERAGE(preco) FROM Produtos;}
        \item \texttt{SELECT AVG(preco) FROM Produtos;}
        \item \texttt{SELECT SUM(preco) FROM Produtos;}
        \item \texttt{SELECT COUNT(preco) FROM Produtos;}
    \end{enumerate}

    \item Com base no modelo lógico apresentado, julgue as afirmativas abaixo sobre a junção de tabelas como Verdadeiras (\textbf{V}) ou Falsas (\textbf{F}):
    \begin{enumerate}[label=( \quad )]
        \item O comando \texttt{SELECT * FROM Produtos INNER JOIN Pedidos;} sem a utilização da cláusula \texttt{ON} ou \texttt{WHERE} para igualar as chaves resultará em um Produto Cartesiano entre as tabelas.
        \item A consulta \texttt{SELECT p.nome\_produto, pe.quantidade FROM Produtos p INNER JOIN Pedidos pe ON p.id\_prod = pe.fk\_prod;} retorna corretamente o nome do produto e a quantidade vendida para cada pedido realizado.
    \end{enumerate}

    \item Com base no modelo lógico apresentado, julgue as afirmativas abaixo sobre funções de agregação e agrupamento como Verdadeiras (\textbf{V}) ou Falsas (\textbf{F}):
    \begin{enumerate}[label=( \quad )]
        \item A cláusula \texttt{WHERE} pode ser utilizada para filtrar o resultado de uma função de agregação como \texttt{AVG()}, desde que esteja posicionada antes do \texttt{GROUP BY}.
        \item A consulta \texttt{SELECT categoria, AVG(preco) FROM Produtos GROUP BY categoria HAVING AVG(preco) > 100;} apresenta a média de preços por categoria, exibindo apenas as que possuem média superior a R\$ 100,00.
    \end{enumerate}

\end{enumerate}

\vspace{0.5cm} 

\end{document} 